\zadanie{Obsadzenia poziomów rotacyjnych HCl}

W jakiej temperaturze liczba cząsteczek HCl w pierwszym rotacyjnym poziomie
wzbudzonym jest równa $1/e$ liczby cząsteczek w rotacyjnym poziomie
podstawowym? Energia poziomu rotacyjnego o liczbie kwantowej $J$ wynosi
\begin{equation*}
  \varepsilon_J = B\,J(J+1), \qquad J = 0, 1, 2, \ldots,
\end{equation*}
gdzie $B = \hbar^2/(2I)$ jest stałą rotacyjną, a $I$ jest momentem
bezwładności cząsteczki. Charakterystyczna temperatura rotacji HCl jest równa
\begin{equation*}
  \theta^\text{rot} = \frac{B}{k_B} = \frac{\hbar^2}{2 I k_B} = 15.02~\text{K} .
\end{equation*}
Prawidłowa odpowiedź: $T = 14.3$~K.

\begin{wskazowki}
\begin{itemize}
  \item Jaka jest degeneracja $g_J$ poziomu $J$? Ile wartości przyjmuje
    liczba kwantowa $M$ dla danego $J$?
  \item W rozkładzie kanonicznym liczba cząsteczek na poziomie $J$ jest
    proporcjonalna do $g_J\,e^{-\varepsilon_J/k_BT}$.
\end{itemize}
\end{wskazowki}
